\documentclass[11pt]{article}

% ============================================================
%  Full faithful rewrite of sPHENIX analysis note PPG-12
%  "Measurement of isolated-prompt photons in p+p collisions
%   at sqrt(s) = 200 GeV with the sPHENIX detector"
%  Internal tag: sPH-ppg-2024-012 (v4)
% ============================================================

\usepackage[margin=1in]{geometry}
\usepackage{amsmath,amssymb}
\usepackage{graphicx}
\usepackage{longtable}
\usepackage{booktabs}
\usepackage{array}
\usepackage{xcolor}
\usepackage[hidelinks]{hyperref}
\usepackage{caption}
\usepackage{enumitem}

\newcommand{\ET}{E_{\mathrm{T}}}
\newcommand{\ETg}{E_{\mathrm{T}}^{\gamma}}
\newcommand{\Eiso}{E_{\mathrm{iso}}}
\newcommand{\pt}{p_{\mathrm{T}}}
\newcommand{\sqrts}{\sqrt{s}}
\newcommand{\dr}{\Delta R}

\title{\textbf{Measurement of Isolated-Prompt Photons in $p+p$ Collisions
at $\sqrts = 200$~GeV with the sPHENIX Detector}\\[4pt]
\large A Faithful, Reader-Friendly Rewrite of the Analysis Note}
\author{Original authors: Yeonju Go, Shuhang Li, Blair Seidlitz,\\
Justin Bennett, Muhammad Elsayed\\[4pt]
\small Internal note tag: sPH-ppg-2024-012 (version 4)}
\date{Original date: 15 May 2026}

\begin{document}
\maketitle

\begin{abstract}
\noindent
This note presents the first measurement of the isolated-prompt-photon
differential cross section in $p+p$ collisions at a center-of-mass energy
$\sqrts = 200$~GeV recorded by the sPHENIX detector at RHIC during Run~2024.
Photons are measured within a pseudorapidity acceptance $|\eta^{\gamma}| < 0.7$
and a transverse-energy range $12 < \ETg < 32$~GeV. The isolation requirement
is defined at truth (generator) level as a total isolation energy
$\Eiso^{\mathrm{truth}} < 4$~GeV inside a cone of radius $\dr = 0.3$ around the
photon. The measurement uses an integrated luminosity of
$\mathcal{L} = 64.37~\mathrm{pb}^{-1}$. The cross section is corrected for
background, detector efficiency and resolution (via iterative Bayesian
unfolding) and compared to next-to-leading-order pQCD predictions from
\textsc{jetphox} and to \textsc{pythia-8}. This rewrite preserves every section,
table and numerical result of the original note while expanding the prose for a
reader with a high-energy-physics graduate background.
\end{abstract}

\tableofcontents
\newpage

% ============================================================
\section{Introduction}
% ============================================================

The production of prompt photons in hadronic collisions is a classic probe of
perturbative quantum chromodynamics (pQCD) and of the parton structure of the
proton. Unlike hadrons, a prompt photon emerges directly from the hard
scattering and does not fragment or hadronize; it therefore carries clean
kinematic information about the underlying partonic process. ``Prompt''
photons are conventionally separated into two contributions:

\begin{itemize}
  \item \textbf{Direct photons}, produced in the hard $2\to2$ scattering
        itself. At RHIC energies the two dominant subprocesses are
        \emph{quark--gluon Compton scattering} $q g \to q \gamma$ and
        \emph{quark--antiquark annihilation} $q\bar q \to g\gamma$. The Compton
        channel involves a gluon in the initial state and therefore makes the
        cross section directly sensitive to the gluon parton distribution
        function (PDF) of the proton.
  \item \textbf{Fragmentation photons}, radiated collinearly from a final-state
        quark or gluon. These are described by photon fragmentation functions
        and are suppressed by requiring the photon to be \emph{isolated}.
\end{itemize}

\noindent
The original note's Fig.~1 shows the representative leading-order Feynman
diagrams for these processes: the Compton diagram ($qg\to q\gamma$), the
annihilation diagram ($q\bar q\to g\gamma$), and a fragmentation diagram in
which a photon is emitted from a scattered parton.

Because the Compton process dominates at RHIC, isolated-photon production probes
the gluon density of the proton. The momentum fraction carried by the struck
partons is approximately
\begin{equation}
  x \;\approx\; \frac{2\,\ETg}{\sqrts},
\end{equation}
so the kinematic range of this measurement, $12 < \ETg < 32$~GeV at
$\sqrts = 200$~GeV, corresponds to $x \approx 0.1$--$0.3$. This is precisely the
region where the gluon PDF is comparatively poorly constrained, making the
measurement valuable as a future input to global PDF fits.

\paragraph{Isolation.} An \emph{isolation} criterion is imposed to suppress the
fragmentation contribution and the much larger background from neutral-meson
decays (primarily $\pi^0 \to \gamma\gamma$ and $\eta\to\gamma\gamma$). One sums
the energy in a cone of fixed radius around the photon, excluding the photon
itself, and requires this isolation energy to be small. The original note's
Figs.~2 and~3 illustrate, from simulation, how the fraction of photons passing
the isolation cut evolves with $\ETg$ and how it differs between signal direct
photons and the meson-decay background.

\paragraph{Prior measurements.} The PHENIX experiment previously measured the
isolated direct-photon cross section in $p+p$ at $\sqrts=200$~GeV
(Phys.\ Rev.\ D \textbf{86}, 072008, Ref.~[1]). sPHENIX, with its hermetic
electromagnetic and hadronic calorimetry, extends the reach to higher $\ETg$
with substantially improved jet/photon energy resolution, and the present note
constitutes the first such sPHENIX result.

% ============================================================
\section{Data and Monte Carlo Samples}
% ============================================================

\subsection{Data sample}

The analysis uses $p+p$ collision data recorded by sPHENIX during the Run~2024
(Run~24) data-taking period. The relevant run range, software production tag and
calibration database (CDB) global tag are summarized in Table~\ref{tab:data}.

\begin{table}[h]
\centering
\caption{Data sample definition.}
\label{tab:data}
\begin{tabular}{ll}
\toprule
Quantity & Value \\
\midrule
Run range & 47289--53880 \\
Analysis build & ana521 \\
Calibration global tag (CDB) & 2025p007 \\
DST datasets & DST\_FITTING, DST\_JETCALO \\
\bottomrule
\end{tabular}
\end{table}

\subsection{Monte Carlo samples}

Signal and background processes are generated with \textsc{pythia-8} (Detroit
tune, version 8.307) and passed through a full \textsc{geant-4} simulation of the
sPHENIX detector. To obtain statistics across the full $\ETg$ spectrum, samples
are generated in separate hard-process bins: photon-triggered samples with
parton-level thresholds of 5, 10 and 20~GeV, and QCD dijet samples with
thresholds of 8, 12, 20, 30 (and 40)~GeV, each with $\mathcal{O}(10^7)$ events.

Each generated sample carries an effective cross section that accounts for the
generator-level acceptance of its hard-process selection,
\begin{equation}
  \sigma_{\mathrm{eff}} \;=\; \sigma_{\mathrm{gen}}\,
  \frac{N_{\mathrm{acc}}}{N_{\mathrm{gen}}},
\end{equation}
where $\sigma_{\mathrm{gen}}$ is the generated process cross section,
$N_{\mathrm{gen}}$ the number of generated events and $N_{\mathrm{acc}}$ the
number accepted into the sample. The effective cross sections are listed in
Table~\ref{tab:sigeff}.

\begin{table}[h]
\centering
\caption{Effective cross sections $\sigma_{\mathrm{eff}}$ [pb] of the MC samples
(MDC2 single-interaction / double-interaction productions).}
\label{tab:sigeff}
\begin{tabular}{lr}
\toprule
Sample & $\sigma_{\mathrm{eff}}$ [pb] \\
\midrule
jet 8  & $4.93\times10^{6}$ \\
jet 12 & $1.49\times10^{6}$ \\
jet 20 & $6.26\times10^{4}$ \\
jet 30 & $2.53\times10^{3}$ \\
photon 5  & $1.46\times10^{5}$ \\
photon 10 & $6.94\times10^{3}$ \\
photon 20 & $1.30\times10^{2}$ \\
\bottomrule
\end{tabular}
\end{table}

To avoid double-counting when stitching the threshold-binned samples together,
each sample is restricted to a truth $\pt$ window (the photon or leading-jet
truth $\pt$), given in Table~\ref{tab:ptwin}.

\begin{table}[h]
\centering
\caption{Truth $\pt$ windows used to stitch the threshold-binned MC samples.}
\label{tab:ptwin}
\begin{tabular}{ll}
\toprule
Sample & Truth $\pt$ window [GeV] \\
\midrule
photon 5  & $[0,14)$ \\
photon 10 & $[14,22)$ \\
photon 20 & $[22,\infty)$ \\
jet 8  & $[9,14)$ \\
jet 12 & $[14,21)$ \\
jet 20 & $[21,32)$ \\
jet 30 & $[32,42)$ \\
jet 40 & $[42,\infty)$ \\
\bottomrule
\end{tabular}
\end{table}

\subsection{Beam crossing angle and double interactions}

The Run~24 data were taken under two distinct beam configurations distinguished
by the crossing angle, which strongly affects the longitudinal width of the
luminous region and hence the rate of double (pileup) interactions:

\begin{itemize}
  \item \textbf{0~mrad crossing angle}: runs 47289--51274,
        $\mathcal{L}\approx 47~\mathrm{pb}^{-1}$, $\sigma_z\approx 54$~cm,
        with a double-interaction (DI) fraction of about 22\% (precisely
        22.4\%).
  \item \textbf{1.5~mrad crossing angle}: runs 51274--54000,
        $\mathcal{L}\approx 17~\mathrm{pb}^{-1}$, $\sigma_z\approx 22$~cm,
        with a DI fraction of about 8\% (precisely 7.9\%).
\end{itemize}

\noindent
Double interactions are modeled in MC by \emph{blending} single-interaction (SI)
and double-interaction templates; this is described in detail in
Sec.~\ref{sec:syst} and Appendix~17.

\subsection{Trigger}

The analysis relies on the EMCal photon (single-tower / cluster) trigger. A Level-1
trigger update was deployed at run~47289; the nominal trigger is the
\textbf{Photon 4~GeV} trigger (trigger bit~30). The per-event trigger efficiency
is measured with a reference (minimum-bias) trigger using the ratio
\begin{equation}
  \varepsilon_{\mathrm{trig}}
  = \frac{N(\text{scaled[10] \& live[}N_{\mathrm{bit}}\text{]})}
         {N(\text{scaled[10]})},
\end{equation}
i.e.\ the fraction of minimum-bias-selected events that also fire the photon
trigger bit. The efficiency turn-on as a function of $\ET$ is fit with a Gumbel
function,
\begin{equation}
  \varepsilon_{\mathrm{trig}}(\ET) = p_0\,
  \exp\!\left[-\exp\!\left(-\frac{\ET-\mu}{\beta}\right)\right],
\end{equation}
reaching a plateau efficiency of $0.9966$ in the measured $\ETg$ range.

\subsection{Event selection and luminosity}

Events are required to have a reconstructed minimum-bias-detector (MBD) vertex
satisfying $|z_{\mathrm{MBD}}| < 60$~cm. The integrated luminosity for the
analyzed sample is
\begin{equation}
  \mathcal{L} = 64.37^{\,+5.88}_{\,-4.34}~\mathrm{pb}^{-1}
  \quad (+9.13\%/-6.75\%),
\end{equation}
determined from a van~der~Meer scan with an MBD reference cross section
$\sigma_{\mathrm{MBD}} = 25.2^{+2.3}_{-1.7}$~mb (see Sec.~\ref{sec:results} and
the PPG-10 luminosity note, Ref.~[10]).

% ============================================================
\section{Photon Reconstruction and Identification}
% ============================================================

\subsection{Calorimeters}

The photon reconstruction relies on the sPHENIX electromagnetic calorimeter
(EMCal) and on the inner and outer hadronic calorimeters (iHCal, oHCal). The
EMCal is a SPACAL-type tungsten/scintillating-fiber sampling calorimeter with full
$2\pi$ azimuthal coverage and pseudorapidity coverage $|\eta|<1.1$. It is
segmented into towers of size $\Delta\eta\times\Delta\phi\approx0.024\times0.024$
($96\times256$ towers in total) and provides a single-photon energy resolution of
$\sigma_E/E\simeq16\%/\sqrt{E[\mathrm{GeV}]}\oplus5\%$ at test-beam energies. As
noted in Sec.~\ref{sec:syst}, an additional smearing is applied to the EMCal
energies in the MC for the nominal analysis to better match the in-situ
resolution. The two HCal sections sit immediately behind the EMCal at increasing
radius; they provide hadronic shower containment and an energy reading that enters
the hadronic-leakage variable $R_{\mathrm{had}}$ and the cluster isolation. The
collision vertex is reconstructed from the Minimum Bias Detector (MBD), two arrays
of $64$ quartz-Cherenkov photomultipliers in the forward regions
$3.51 < |\eta| < 4.61$.

\subsection{Clustering algorithm}

EMCal towers are aggregated into clusters by the
\texttt{RawClusterBuilderTemplate} algorithm. Edge-adjacent towers above a single
$70$~MeV threshold that pass the \texttt{isGood} tower-quality cuts are grouped
together; clusters grow until no adjacent tower exceeds the threshold, after which
a cluster-splitting algorithm separates overlapping showers. The cluster
kinematics use the tower positions in $x$--$y$--$z$ space with shower-depth
corrections, and the resulting position is converted to $\eta$ and $\phi$ relative
to the MBD-reconstructed collision vertex; the cluster energy is the scalar sum of
the constituent tower energies. The dependence of the cluster $\eta$ on the vertex
$z$ position is the mechanism by which the double-interaction vertex shift
(Appendix~\ref{app:di}) distorts the reconstructed shower shapes.

\subsection{Photon energy response}

The energy scale and resolution are estimated from MC. The ratio of reconstructed
cluster $\ET$ to truth photon $\ET^{\mathrm{truth}}$ is studied as a function of
$\ET^{\mathrm{truth}}$ and fit with a Gaussian in each truth-energy bin (original
note's Figs.~8--10). To isolate the intrinsic detector resolution from
vertex-induced distortions, this study restricts the vertex to $|z_{\mathrm{vtx}}|<30$~cm;
the wider analysis fiducial $|z_{\mathrm{vtx}}|<60$~cm is restored in the unfolding
response matrix and the efficiency calculation. The fitted mean is close to unity
($\mu\approx0.99$) and the fractional resolution $\sigma$ improves from about
$0.071$ at $8$--$10$~GeV to about $0.051$ at $36$--$45$~GeV. A low-side tail
reflects residual position-dependent calibration and is absorbed into the
unfolding response matrix rather than corrected separately. The
signal-to-background ratio of the reconstructed sample rises from $\sim15\%$ near
$10$~GeV to $\sim50\%$ near $30$~GeV, which is the underlying reason the extracted
purity (Sec.~\ref{sec:analysis}) improves so strongly with $\ETg$.

\subsection{Shower-shape variables}

A set of EMCal shower-shape observables is used to distinguish single photons
from the broader and more irregular showers produced by merged
$\pi^0\to\gamma\gamma$ decays and by hadronic background. These variables are
summarized in Table~\ref{tab:showershape}.

\begin{table}[h]
\centering
\caption{Shower-shape variables used in photon identification.}
\label{tab:showershape}
\begin{tabular}{ll}
\toprule
Variable & Description \\
\midrule
$E_{3\times2}/E_{3\times5}$ & ratio of energy in $3\times2$ to $3\times5$ tower windows \\
$E_{1\times1}/E$            & central-tower energy fraction \\
$w_{\mathrm{COGX}\,\eta}$   & shower width (center-of-gravity) in $\eta$ \\
$w_{\mathrm{COGX}\,\phi}$   & shower width (center-of-gravity) in $\phi$ \\
$w_{72}$                    & 7$\times$2 lateral width \\
$R_{\mathrm{had}}$          & hadronic-leakage ratio (HCal behind cluster) \\
$E_1$--$E_4$                & ring energies around the cluster core \\
$E_{t1}$--$E_{t4}$          & corresponding transverse ring energies \\
BDT                         & combined boosted-decision-tree discriminant \\
\bottomrule
\end{tabular}
\end{table}

\subsection{Non-physics-background (NPB) rejection}

A first XGBoost boosted-decision-tree (BDT) discriminant removes
\emph{non-physics background} (NPB): out-of-time energy depositions, noise and
beam-induced backgrounds. The NPB BDT uses 25 input features and is trained
using the cluster timing as a label: clusters with
$t_{\mathrm{cluster}} - t_{\mathrm{MBD}} < -7$~ns are tagged as NPB. Training and
evaluation use $6 < \ET < 40$~GeV and $|\eta| < 0.7$. The nominal working point
is $\mathrm{NPB} > 0.5$, which gives roughly $99\%$ purity at roughly $99\%$
signal retention; the cut is varied to $0.3$ and $0.7$ for systematics.

\subsection{Preselection and photon-ID BDT}

After NPB rejection, photons pass a loose \emph{preselection} and are then
classified by a dedicated photon-ID BDT. The preselection and the
tight/non-tight BDT working-point definitions are given in
Table~\ref{tab:presel}. The non-tight selection is used as the control sample in
the data-driven purity (ABCD) method (Sec.~\ref{sec:analysis}).

\begin{table}[h]
\centering
\caption{Preselection requirements and tight / non-tight photon-ID working
points. The tight and non-tight thresholds are linear in $\ETg$.}
\label{tab:presel}
\begin{tabular}{ll}
\toprule
Selection & Requirement \\
\midrule
\multicolumn{2}{l}{\textit{Preselection}}\\
\quad central fraction & $E_{1\times1}/E_{3\times3} < 0.98$ \\
\quad ring 1 transverse & $0.6 < e_{t1} < 1.0$ \\
\quad window ratio & $0.8 < E_{3\times2}/E_{3\times5} < 1.0$ \\
\quad NPB BDT & $\mathrm{NPB} > 0.5$ \\
\midrule
\multicolumn{2}{l}{\textit{Tight}}\\
\quad photon-ID BDT & $\mathrm{BDT} > 0.8156 - 0.00156\,\ETg$ \\
\midrule
\multicolumn{2}{l}{\textit{Non-tight}}\\
\quad photon-ID BDT & $0.7333 - 0.01333\,\ETg < \mathrm{BDT} < 0.6844 + 0.00156\,\ETg$ \\
\bottomrule
\end{tabular}
\end{table}

The photon-ID BDT (Table~\ref{tab:idbdt}) uses 11 input features and consists of
approximately 750 trees of depth~5.

\begin{table}[h]
\centering
\caption{Input features of the photon-ID BDT (11 features).}
\label{tab:idbdt}
\begin{tabular}{ll}
\toprule
\# & Feature \\
\midrule
1  & $\ET$ \\
2  & $w_{\mathrm{COGX}\,\eta}$ \\
3  & $w_{\mathrm{COGX}\,\phi}$ \\
4  & $z_{\mathrm{vtx}}$ \\
5  & $\eta$ \\
6  & $E_{1\times1}/E_{3\times3}$ \\
7--10 & $E_{t1}$--$E_{t4}$ \\
11 & $E_{3\times2}/E_{3\times5}$ \\
\bottomrule
\end{tabular}
\end{table}

\subsection{Isolation}

Several isolation definitions were evaluated by their ROC area-under-curve
(AUC); the best performer is a TopoCluster isolation in a cone of radius
$R=0.4$, with AUC in the range $\sim0.82$--$0.88$. The nominal $80\%$-efficiency
isolation working point is
\begin{equation}
  \Eiso < 0.490 + 0.0370\,E_{\mathrm{reco}}\;\mathrm{GeV},
\end{equation}
and the complementary non-isolated control region is
\begin{equation}
  \Eiso > 0.8 + 0.490 + 0.0370\,E_{\mathrm{reco}}\;\mathrm{GeV}.
\end{equation}
The isolation cone is referenced to the EMCal radius of $93.5$~cm.


% ============================================================
\section{Analysis Procedure}
\label{sec:analysis}
% ============================================================

The measurement proceeds through an eight-step workflow: (1) event and photon
selection; (2) NPB rejection; (3) photon-ID and isolation classification into
ABCD regions; (4) data-driven purity extraction with signal-leakage correction;
(5) MC closure of the purity method; (6) efficiency determination
(reconstruction, ID, isolation, MBD-vertex); (7) iterative Bayesian unfolding;
and (8) normalization to the differential cross section.

\subsection{Truth isolation}

At generator level a photon is defined as isolated if the total isolation energy
inside a cone $\dr = 0.3$ around the photon satisfies
$\Eiso^{\mathrm{truth}} < 4$~GeV. This truth definition is the reference to which
the unfolded result is corrected, decoupling the published cross section from
the detector-level reconstruction (TopoCluster $R=0.4$) isolation.

\subsection{Data-driven purity: the double-sideband (ABCD) method}

The signal purity is extracted with a two-dimensional double-sideband (``ABCD'')
method using the photon-ID (tight/non-tight) and isolation (isolated/non-isolated)
axes. The four regions are:
\begin{itemize}[noitemsep]
  \item \textbf{A}: tight \& isolated (signal region),
  \item \textbf{B}: tight \& non-isolated,
  \item \textbf{C}: non-tight \& isolated,
  \item \textbf{D}: non-tight \& non-isolated.
\end{itemize}
Assuming the ID and isolation variables are uncorrelated for background, the
background in region~A is estimated from the sidebands, giving the signal yield
\begin{equation}
  N_A^{\mathrm{sig}} = N_A - N_B\,\frac{N_C}{N_D}.
\end{equation}
A residual signal-leakage correction accounts for genuine signal photons leaking
into the control regions (estimated from MC fractions $f_X^{\mathrm{MC}}$, with
$X\in\{B,C,D\}$):
\begin{equation}
  N_A^{\mathrm{sig}} = \frac{N_A - N_B\,\dfrac{N_C}{N_D}}
  {1 - \big(f_B - f_D\tfrac{N_C}{N_D}\big) - \dots},
\end{equation}
(schematically; the full leakage-corrected expression uses the MC signal
fractions in each region). The purity is then
\begin{equation}
  P = \frac{N_A^{\mathrm{sig}}}{N_A}.
\end{equation}

The extracted purity rises strongly with $\ETg$:
\begin{itemize}[noitemsep]
  \item $P = 0.51 \pm 0.01$ at $10$--$12$~GeV,
  \item $P = 0.72 \pm 0.04$ at $18$--$20$~GeV,
  \item $P = 0.90 \pm 0.05$ above $26$~GeV.
\end{itemize}

\subsection{MC closure and isolation-template tuning}

The ABCD procedure is validated by a closure test in MC, $C_{\mathrm{MC}}$, in
which the method is applied to simulation and the recovered purity is compared to
the true signal fraction. To make the simulated isolation distribution match
data, the MC isolation energy is corrected with a linear ``fudge'' transformation
\begin{equation}
  \Eiso^{\mathrm{MC,corr}} = 1.2\,\Eiso^{\mathrm{MC}} + 0.1~\mathrm{GeV}.
\end{equation}

\subsection{Efficiencies}

Four efficiency factors are evaluated from MC as functions of the truth photon
$\ETg$ and applied bin-by-bin (together with the purity) to correct the measured
yields. Truth photons are matched to reconstructed clusters using the
track-ID-based association from the \texttt{CaloRawClusterEval} module, which
assigns to each cluster the truth primary particle whose \textsc{geant} hits
contribute the largest summed energy across the cluster's towers; a truth photon
is matched when its track ID equals the cluster's stored leading-energy truth ID.
The four factors are defined as follows.
\begin{itemize}[noitemsep]
  \item \textbf{Reconstruction efficiency}: the fraction of truth signal photons
        matched to a reconstructed cluster.
  \item \textbf{Identification efficiency}: the fraction of truth-matched signal
        photons that also pass the preselection and the tight photon-ID, relative
        to all truth-matched signal photons.
  \item \textbf{Isolation efficiency}: the fraction of truth-matched signal
        photons that also satisfy the isolation requirement, relative to all
        truth-matched signal photons.
  \item \textbf{MBD-vertex efficiency}: the fraction of truth-isolated signal
        photon events in which the MBD north--south coincidence fires and a
        reconstructed vertex is found within $|z_{\mathrm{reco}}|<60$~cm. The
        denominator uses the all-$z$ truth fiducial, so this factor converts the
        beam-delivered luminosity into events with a usable reconstructed vertex.
\end{itemize}
The reconstruction, ID and isolation efficiencies are each in the $0.7$--$0.9$
range and vary slowly with $\ETg$; the MBD-vertex efficiency is the factor
studied for channel independence in Appendix~\ref{app:mbd}. The efficiencies are
applied after unfolding, on the truth-$\ETg$ grid.

\subsection{Unfolding}

The detector-level spectrum is corrected back to truth level with the D'Agostini
iterative Bayesian method as implemented in RooUnfold~3.0.5. The iterative method
is preferred over matrix inversion or SVD for two reasons: it converges rapidly
when the response matrix is close to diagonal --- as it is here, given the
$\sim5\%$ EMCal energy resolution --- and the regularization is set transparently
by the iteration count, with each additional iteration trading residual bias
against statistical noise. The fiducial binning of the reported result is
\begin{equation}
  [12,14,16,18,20,22,24,26,28,32]~\mathrm{GeV}.
\end{equation}
The response matrix is built on a wider grid that adds a $[10,12]$~GeV reco-side
underflow and a $[32,36]$~GeV reco-side overflow, paired with truth bins
$[8,10]$ (low underflow) and $[36,45]$~GeV (high overflow); these absorb migration
in and out of the reported range and are not themselves reported.

To make the MC $\ETg$ prior resemble data, the response matrix is reweighted by
the ratio of a single-iteration D'Agostini unfold of the data spectrum (from a
baseline run with the prior reweight disabled) to the un-reweighted MC truth
prior, both normalized to unit area. The ratio is fit with an order-$[2/2]$
Pad\'e approximant from $10$ to $35$~GeV (with the rational denominator's zeros
constrained below the $[10,30]$~GeV clamp range), and the factor is applied
photon-by-photon on the truth $\ETg$ when populating the response matrix in
the next pipeline iteration. Because the energy resolution is small, the effect of
unfolding on the spectrum is minor compared with the statistical uncertainties:
the relative change between successive iterations drops below the per-bin
statistical uncertainty after the first iteration, so the nominal result uses
$N_{\mathrm{iter}}=2$, which keeps the iteration-dependent bias subdominant to
other systematics. The ratio of post- to pre-unfolding yield is $2$--$10\%$ across
the $\ETg$ range.

Two closure tests validate the procedure. In the \emph{full} closure test the
response matrix is built from the entire MC sample and the same events are
unfolded, exactly reproducing the truth spectrum. In the \emph{half} closure test
the matrix is built from half the sample and applied to the statistically
independent other half; the unfolded result agrees with truth to within $2\%$,
consistent with the statistical uncertainty. The iteration scan that sets
$N_{\mathrm{iter}}$ is propagated as the unfolding-iteration systematic
(Sec.~\ref{sec:syst}, item~5.7) and the prior reweighting as item~5.6.

% ============================================================
\section{Systematic Uncertainties}
\label{sec:syst}
% ============================================================

The systematic uncertainties are grouped by source. Maximum per-bin values
(quoted as down/up where asymmetric) are collected in Table~\ref{tab:syst}.

\begin{description}
  \item[5.1 Energy scale (flat).] A $\pm1.1\%$ shift of the EMCal energy scale
        produces a maximum cross-section variation of $-10.5\%/+12.2\%$.
  \item[5.1 Energy scale (non-linearity).] A non-linearity slope $s=5\times10^{-4}$
        gives up to $\pm7.2\%$.
  \item[5.2 Energy resolution.] The nominal MC smearing
        $\sigma_{\mathrm{MC}}=(0.185,0,0.040)$ is compared to a data-like
        $\sigma_{\mathrm{data}}=(0.15,0.05,0.05)$, with variations of no extra
        smearing and $(0.13,0.08,0.08)$; the maximum effect is $\pm9.6\%$
        (up to $\pm21.2\%$ in the most extreme bin).
  \item[5.3.1 Tight-ID variation.] Shifting the tight BDT threshold by $+0.05$
        gives $\pm10.4\%$.
  \item[5.3.2 Non-tight variation.] Shifting the non-tight threshold by $-0.10$
        gives $\pm2.1\%$.
  \item[5.3.3 Non-isolation definition.] Using \texttt{noniso04}/\texttt{noniso10}
        control regions gives $-5.1\%/+0.9\%$.
  \item[5.3.4 Fit functional form.] Using an error-function fit form gives
        $\pm0.7\%$.
  \item[5.3.5 Fit confidence interval.] The purity-fit CI gives $-8.1\%/+8.2\%$.
  \item[5.3.6 MC closure.] The purity-method closure gives $\pm12.0\%$ (largest
        at low $\ETg$).
  \item[5.4.1 Isolation pedestal.] $\pm5.0\%$.
  \item[5.4.2 \textsc{herwig} cross-check.] Re-deriving efficiencies with
        \textsc{herwig}~7 gives $\pm5.3\%$.
  \item[5.5 MBD vertex.] $\pm6\%$.
  \item[5.6 Response reweighting.] $\pm0.8\%$.
  \item[5.7 Unfolding iteration scan.] $\pm4.7\%$.
  \item[5.8 Double-interaction blending.] Best-fit blend fractions
        $f_{\mathrm{best}}=0.290$ (0~mrad) and $0.000$ (1.5~mrad) give $\pm1.6\%$.
  \item[5.9 NPB cut.] Varying the NPB working point gives $-1.4\%/+6.1\%$.
  \item[5.10 Luminosity.] $+9.13\%/-6.75\%$.
\end{description}

\begin{longtable}{lcc}
\caption{Maximum per-bin systematic uncertainties [\%] (down / up where
applicable).}\label{tab:syst}\\
\toprule
Source & Max down [\%] & Max up [\%] \\
\midrule
\endfirsthead
\toprule
Source & Max down [\%] & Max up [\%] \\
\midrule
\endhead
\bottomrule
\endfoot
Energy scale (flat)        & 10.49 & 11.87 \\
Energy scale (non-lin.)    & 7.26  & 7.26  \\
Energy resolution          & 21.21 & 21.21 \\
Purity -- tight            & \multicolumn{2}{c}{6.72} \\
Purity -- non-tight        & \multicolumn{2}{c}{2.35} \\
Purity -- non-iso          & 3.80  & 0.99  \\
Purity -- fit form         & \multicolumn{2}{c}{0.33} \\
Purity -- fit stat (CI)    & \multicolumn{2}{c}{4.40} \\
Purity -- MC closure       & \multicolumn{2}{c}{5.77} \\
Efficiency -- isolation    & \multicolumn{2}{c}{4.73} \\
Efficiency -- \textsc{herwig} & \multicolumn{2}{c}{5.83} \\
MBD vertex                 & \multicolumn{2}{c}{6.00} \\
Unfolding -- reweight      & \multicolumn{2}{c}{0.36} \\
Unfolding -- iteration     & \multicolumn{2}{c}{5.50} \\
Double interaction         & \multicolumn{2}{c}{1.57} \\
NPB cut                    & 1.03  & 4.54  \\
Luminosity                 & 6.75  & 9.13  \\
\midrule
\textbf{Total}             & \textbf{30.80} & \textbf{31.72} \\
\end{longtable}

% ============================================================
\section{Results}
\label{sec:results}
% ============================================================

The reconstructed yield is obtained by unfolding the background-subtracted,
purity-weighted signal,
\begin{equation}
  Y_{\mathrm{rec}} = \mathrm{Unfold}\big(N_{\mathrm{tight,iso}}\cdot P\big),
\end{equation}
and the isolated-prompt-photon differential cross section is
\begin{equation}
  \frac{d\sigma}{d\ETg\,d\eta}
  = \frac{1}{\mathcal{L}}\,
    \frac{Y_{\mathrm{rec}}}{\varepsilon\,\Delta\ETg\,\Delta\eta^{\gamma}},
\end{equation}
where $\varepsilon$ is the combined efficiency, $\Delta\ETg$ the bin width and
$\Delta\eta^{\gamma}=1.4$ the pseudorapidity acceptance ($|\eta^{\gamma}|<0.7$).

The fully corrected differential cross section (original Fig.~54) falls by
roughly two orders of magnitude across the reported $12<\ETg<32$~GeV range. It is
compared to two theory calculations. \textsc{pythia-8} (version~8.307, Detroit
tune) is run with the same truth-level isolation as the data.
\textsc{jetphox} (v1.3.1.4) computes the direct$+$fragmentation cross section at
NLO with CT14nlo PDFs, BFG~II fragmentation functions, and scales
$\mu_R=\mu_F=\mu_f=\ETg$; an additional NLO prediction without isolation is
overlaid for reference. Both calculations agree with the measurement within
experimental uncertainties across the full $\ETg$ range, supporting the
gluon-Compton-dominance picture at the parton momentum fractions probed here.

The measurement is systematics-dominated over the whole range. The statistical
uncertainty grows from sub-percent at the lowest reported bin to about $30\%$ at
the highest but remains subdominant. The energy-scale and double-interaction
systematics dominate the high-$\ETg$ budget, while the purity (fit-CI and
MC-closure) and luminosity sources drive the low-$\ETg$ budget; the total
systematic ranges from about $\pm14\%$ at low $\ETg$ to $-30.8\%/+31.7\%$ near the
top of the range (Table~\ref{tab:syst}).

The result is also compared to the earlier PHENIX measurement (Ref.~[1], original
Fig.~55), which reports inclusive direct photons (no isolation) in $|\eta|<0.25$.
The PHENIX spectrum is shown both as published and after the bin-width recovery
and pseudorapidity-density rescaling of Appendix~\ref{app:phenix}; the corrected
spectrum is consistent with the present measurement within uncertainties across
the overlap range, with the small residual offset attributable to the
fragmentation contribution that survives the truth-level isolation cut.

% ============================================================
\section{Conclusion}
% ============================================================

This note presents the first sPHENIX measurement of the isolated-prompt-photon
differential cross section in $p+p$ collisions at $\sqrts=200$~GeV, for
$|\eta^{\gamma}|<0.7$ and $12<\ETg<32$~GeV, with a truth isolation
$\Eiso^{\mathrm{truth}}<4$~GeV in $\dr=0.3$ and an integrated luminosity of
$64.37~\mathrm{pb}^{-1}$. The result is corrected for background via a
data-driven double-sideband purity method, for efficiency, and for detector
resolution via iterative Bayesian unfolding, and is found to be consistent with
NLO pQCD (\textsc{jetphox}) and with \textsc{pythia-8}, as well as with the
earlier PHENIX measurement. The measurement demonstrates the photon
reconstruction and identification capabilities of the sPHENIX calorimeter system
and provides a foundation for future gluon-PDF-sensitive measurements at RHIC.

% ============================================================
\section*{References}
\addcontentsline{toc}{section}{References}
% ============================================================
\begin{enumerate}[label={[\arabic*]}]
  \item PHENIX Collaboration, isolated direct photons in $p+p$ at
        $\sqrts=200$~GeV, Phys.\ Rev.\ D \textbf{86}, 072008.
  \item \textsc{pythia} 8.2, an introduction to PYTHIA 8.2.
  \item \textsc{geant4}, a simulation toolkit.
  \item M.\ Spousta and B.\ Cole, jet-substructure / photon studies.
  \item sPHENIX EMCal test-beam results.
  \item ATLAS topo-cluster reconstruction, Eur.\ Phys.\ J.\ C \textbf{77}, 490.
  \item ATLAS isolated-photon cross section, JHEP \textbf{10} (2019) 203.
  \item G.\ D'Agostini, iterative Bayesian unfolding.
  \item T.\ Adye, RooUnfold.
  \item sPHENIX PPG-10 luminosity note.
  \item ATLAS electron/photon performance, JINST \textbf{14}, P12006.
  \item Detroit tune, Phys.\ Rev.\ D \textbf{105}, 016011.
  \item BFG photon fragmentation functions.
\end{enumerate}


\newpage
% ============================================================
\appendix
\section*{Appendices}
\addcontentsline{toc}{section}{Appendices}
% ============================================================

\section{NPB purity and efficiency scans}
\label{app:npb}
The non-physics-background (NPB) discriminant introduced in
Sec.~\ref{sec:analysis} is a dedicated XGBoost classifier whose purpose is to
remove energy depositions in the EMCal that are not associated with a genuine
$p+p$ collision: out-of-time clusters, electronic noise, and beam-induced
``streak'' clusters that align in a narrow band of constant $\eta$ or $\phi$.
These clusters are pathological in that they carry no recoil activity and have
shower shapes that lie outside the physical phase space populated by single
photons or by hadronic showers. Because they survive into the photon-candidate
sample if left untreated, they would inflate the measured yield and bias the
purity extraction; the NPB BDT is therefore applied \emph{before} the photon-ID
BDT and the isolation classification.

The BDT is trained using cluster timing as the supervisory label. The cluster
time is referenced to the MBD event time, and clusters with
$t_{\mathrm{cluster}} - t_{\mathrm{MBD}} < -7$~ns are tagged as NPB, since
genuine collision clusters are tightly peaked in time around the beam crossing
whereas out-of-time and beam-induced deposits populate the early tail. Twenty-five
input features are used, and training and evaluation are restricted to the
fiducial region $6 < \ET < 40$~GeV and $|\eta| < 0.7$, slightly wider in $\ET$
than the analysis window so that the turn-on and the high-$\ET$ tail are both
well sampled.

The working-point study scans the NPB BDT threshold and records, as a function
of $\ETg$, both the photon purity (the fraction of selected clusters that are
genuine physics clusters, estimated from the timing-labeled control sample) and
the signal-retention efficiency (the fraction of genuine clusters that survive
the cut). The nominal cut $\mathrm{NPB}>0.5$ achieves approximately $99\%$
purity while retaining approximately $99\%$ of genuine clusters, an unusually
favorable operating point that reflects the cleanly separated, almost
non-overlapping NPB and physics populations in the timing-trained feature space.
To bracket the residual modeling dependence of this choice the cut is loosened to
$0.3$ and tightened to $0.7$; the resulting per-bin variation of the unfolded
cross section is propagated as the NPB systematic (Sec.~\ref{sec:syst},
item~5.9), where it is found to be small except at the lowest $\ETg$ bin. An
independent, physics-based cross-check of the same BDT using a back-to-back
recoil-jet requirement is presented in Appendix~\ref{app:b2b}.

\section{Shower-shape distributions and correlations}
\label{app:showershape}
This appendix collects the detector-level shower-shape distributions that
motivate the variable choices of Tables~\ref{tab:showershape} and~\ref{tab:idbdt}
and that underpin both the NPB rejection and the photon-ID BDT. The original
note's Figs.~57, 13, 58, 19, 59 and~60 show these distributions for data, signal
MC, inclusive (background-dominated) MC, and an NPB-enriched data template,
together with their pairwise correlations.

\paragraph{Distributions without preselection.} At low $\ETg$ ($14$--$18$~GeV,
original Fig.~57) the shower-shape variables are shown before any preselection
for four overlaid samples: data (black), signal MC (red), inclusive MC (blue),
and an NPB-tagged data template (green). The NPB-tagged template is the subset of
data selected as streak-like by a high center-of-gravity width
$w_{\mathrm{COGX}\,\eta}$ tag; it is normalized to the bulk data by tail-matching
in the region $w_{\mathrm{COGX}\,\eta} > 1$, where the streak population
dominates, and serves as a background-enhanced template for the non-physics
component. The comparison demonstrates that the streak clusters occupy the broad
high-$w$ tail that is essentially absent in the signal MC, which is precisely
what the NPB BDT exploits.

\paragraph{Effect of preselection.} The same distributions are shown after the
loose preselection (original Fig.~58, here at $10$--$14$~GeV). The preselection
suppresses the high-tail population in the variables most sensitive to
non-physics backgrounds, $w_{\mathrm{COGX}\,\eta}$ in particular, while leaving
the bulk of the physics-cluster shapes essentially unchanged. The per-panel
$\chi^2/\mathrm{ndf}$ between inclusive MC and data remains large for several
variables, a reminder that the inclusive MC is a mixture whose composition
(signal fraction, NPB fraction, pileup fraction) is not identical to data; the
purity method (Sec.~\ref{sec:analysis}) is designed to be robust against exactly
this kind of template imperfection, since it extracts the signal yield from the
data sidebands rather than from an absolute MC normalization.

\paragraph{Correlation with isolation.} It is important that the shower-shape
discriminant and the isolation variable be as uncorrelated as possible for
background, because the double-sideband (ABCD) purity method assumes
factorization of the two axes. The original Figs.~59 and~60 show the mean
isolation energy $\langle \Eiso\rangle$ as a function of each shower-shape value,
for background MC, at low ($10$--$14$~GeV) and high ($28$--$30$~GeV) $\ETg$. The
measured linear correlation coefficients are modest, ranging from about $-0.26$
to $+0.42$ depending on the variable and the $\ETg$ range; the width variables
$w_{\mathrm{COGX}\,\eta}$ and $w_{\mathrm{COGX}\,\phi}$ are the most strongly
(positively) correlated with isolation, while the ring transverse energies
$E_{t2}$--$E_{t4}$ are nearly uncorrelated. These residual correlations are the
physical origin of the signal-leakage correction applied to the ABCD result and
are quantified by the MC closure test described in Sec.~\ref{sec:analysis}.

\section{Isolation distributions}
\label{app:isodist}
The reconstructed isolation-energy distribution $\Eiso$ is the second axis of the
double-sideband purity method, and its behavior as a function of $\ETg$ and of
photon-ID category is what makes the ABCD construction work. The original note's
Fig.~61 shows $\Eiso$ in six $\ETg$ bins spanning $10$--$30$~GeV, overlaying three
samples: tight-ID photons in inclusive MC (black), non-tight-ID photons in
inclusive MC (red), and tight-ID photons in signal MC (blue).

Several features support the analysis design. First, the tight signal-MC
distribution is sharply peaked near $\Eiso \approx 0$ with a narrow tail,
confirming that genuine isolated prompt photons deposit very little additional
energy in the cone; this is what justifies the $80\%$-efficiency isolation
working point $\Eiso < 0.490 + 0.0370\,E_{\mathrm{reco}}$~GeV. Second, the
non-tight inclusive sample is shifted to substantially larger $\Eiso$, reflecting
the jet-like environment of the meson-decay and fragmentation backgrounds that
populate the non-tight ID region; this is the basis for using the non-tight
sample as a background-shape control. Third, the separation between the signal
peak and the background tail grows with $\ETg$, which is the microscopic reason
the extracted purity rises from $\sim0.5$ at $10$--$12$~GeV to $\sim0.9$ above
$26$~GeV.

The simulated $\Eiso$ distribution does not match data perfectly, because the
underlying-event activity and the soft-particle spectrum entering the cone are
imperfectly modeled. This residual mismatch is corrected with the linear
``fudge'' transformation
$\Eiso^{\mathrm{MC,corr}} = 1.2\,\Eiso^{\mathrm{MC}} + 0.1$~GeV introduced in
Sec.~\ref{sec:analysis}; the slope and offset are tuned so that the
fudge-corrected MC reproduces the data isolation shape in a background-dominated
control region. The adequacy of this correction is one of the inputs to the
purity MC-closure systematic (Sec.~\ref{sec:syst}, item~5.3.6), and the
sensitivity to the precise pedestal value is captured by the isolation-pedestal
systematic (item~5.4.1).

\section{Shower-shape cut comparison}
\label{app:cutcomp}
This appendix documents the effect of the loose preselection chain on the
shower-shape distributions and, by extension, justifies the choice of a BDT
photon identifier over a traditional sequential (rectangular) cut selection. The
original note's Fig.~62 ($10<\ETg<12$~GeV) and Fig.~63 ($26<\ETg<35$~GeV) show,
for the inclusive-MC background sample, each shower-shape variable before any cut
(blue) and after the full preselection chain of Table~\ref{tab:presel}
(red): the central-fraction requirement $E_{1\times1}/E_{3\times3}<0.98$, the
first-ring transverse energy window $0.6<e_{t1}<1.0$, the window ratio
$0.8<E_{3\times2}/E_{3\times5}<1.0$, and the NPB score requirement
$\mathrm{NPB}>0.5$.

The qualitative message is that the preselection suppresses the high tails of the
variables that are sensitive to non-physics background, most visibly the
center-of-gravity width $w_{\mathrm{COGX}\,\eta}$, while having little effect on
the core of the physics-cluster shapes. In other words, the preselection removes
the pathological clusters without sculpting the discriminating variables of the
genuine population, which is the desired behavior before the multivariate
photon-ID step.

The broader motivation for using a BDT rather than a sequential cut is that the
shower-shape variables are correlated, so a product of one-dimensional
rectangular windows discards discriminating information contained in their joint
distribution. The photon-ID BDT (Table~\ref{tab:idbdt}) combines the eleven
inputs into a single discriminant and, evaluated at fixed signal efficiency,
yields a higher background rejection than the equivalent rectangular selection.
The BDT also allows the working point to depend smoothly on $\ETg$ (the
tight/non-tight thresholds of Table~\ref{tab:presel} are linear in $\ETg$), which
accommodates the $\ETg$-dependence of both the signal shower shapes and the
background composition in a way that fixed cuts cannot.

\section{$\phi$-symmetry tower acceptance}
\label{app:phisym}
This appendix studies the residual data--MC mismatch in the EMCal dead-tower
acceptance and derives a purely data-driven systematic on the cross section from
it. The sPHENIX EMCal contains dead and malfunctioning towers whose response is
not perfectly emulated by the \textsc{geant-4} simulation. If the set of dead
towers active in Run~24 data differs from the set modeled in MC, the
MC-derived reconstruction efficiency is biased high in the affected regions,
which in turn biases the measured cross section low. The $\phi$-symmetry method
exploits the fact that the inclusive cluster yield is $\phi$-symmetric at fixed
$\eta$ in $p+p$ collisions to derive a dead-tower mask directly from data,
without relying on any particular MC dead-tower list. Because the same mask is
then applied identically to numerator and denominator on both the data and MC
sides, the efficiency-corrected cross section becomes sensitive only to the
residual data-versus-MC asymmetry in the masked region, not to the absolute mask
coverage.

\paragraph{Method.} The EMCal is binned in tower index as
$96\,(i_\eta)\times256\,(i_\phi)$. The towers are grouped into $48$ two-row
$\eta$ bands. For each band the $256$ per-$\phi$ cluster counts are projected out.
Under the assumption of $\phi$-symmetry, deviations from a flat distribution
trace either statistical fluctuations or local detector effects. An iterative
$3\sigma$-clipped Gaussian is fit to the $256$-bin distribution, yielding a
band-by-band mean $\mu_{\phi,\mathrm{band}}$ and width $\sigma_{\phi,\mathrm{band}}$.
A tower $(i_\eta,i_\phi)$ is flagged as data-deficient when
$(n_{\mathrm{data}}-\mu_{\phi,b})/\sigma_{\phi,b} < -2$; both $\eta$ rows in the
band are masked whenever a $\phi$ bin is flagged. The original note's Fig.~64
shows an example for the band $i_\eta=64$--$67$ ($\eta\approx0.37$--$0.46$), in
which a contiguous block of $\phi$ cells around $i_\phi\approx30$--$45$ falls well
below the $\mu-2\sigma$ threshold and is tagged.

Four masks are built at successively tighter selection levels: (i)~\emph{preselect}
($\ETg\ge8$~GeV with no-saturation and no shape cuts); (ii)~\emph{common}
(preselect intersected with the NPB cut, $E_{1\times1}/E_{3\times3}<0.98$, and the
width windows); (iii)~\emph{tight} (common intersected with the ten shower-shape
windows and the parametric BDT requirement); and (iv)~\emph{OR} (the union of the
common and tight masks). The corresponding masked-tower counts are $1344$, $1160$,
$646$ and $1180$. The monotonic decrease with selection tightness is the expected
statistical behavior: the tighter samples have lower per-$\phi$ statistics, so the
fit width $\sigma$ is broader and fewer bins fall below the $-2\sigma$ threshold.
The tight mask is largely contained in the common mask (a $626$-tower overlap), so
the OR mask adds only about $20$ towers relative to common alone.

\paragraph{Results.} The mask is applied as a cluster-level veto, identically to
data and MC, and the full efficiency, merge, and yield pipeline is re-run for each
variant. The resulting integrated cross-section shifts, evaluated in
$\ETg\in[10,36]$~GeV relative to the mask-free nominal
$\sigma_{\mathrm{nom}}=1963$~pb, are collected in Table~\ref{tab:phisym}. All four
shifts are positive, as expected: masking removes MC content in data-dead regions
that was inflating the MC-derived efficiency, so the corrected cross section rises
once the asymmetry is removed. The four variants agree to within about $1\%$,
consistent with a uniform detector-level dead-tower effect, and the per-$\ETg$-bin
ratios lie within $\pm3\%$ of unity across $\ETg\in[10,30]$~GeV (original Fig.~67).
The larger scatter in the highest $[32,36)$~GeV bin is statistical and does not
affect the integrated result, which is dominated by the low-$\ETg$ region.

An independent estimator rebins the data and lumi-weighted MC tower maps into
$4\times4$ super-cells and flags super-cells whose $\log(R)$ of the data/MC ratio
deviates below a Gaussian fit to the bulk; this tags $2$--$3\%$ of super-cells per
selection level and yields an integrated shift agreeing with the $\phi$-symmetry
mask to about $1\%$. The convergence of two independent data-driven estimators on
the same $\sim2\%$ level provides a robust bound that does not rely on any specific
MC dead-tower list. The per-bin envelope adopted for the systematic, taken from
the OR mask, is $\pm1.57\%$.

\begin{table}[h]
\centering
\caption{$\phi$-symmetry tower-acceptance cross-check. Nominal cross section
$\sigma_{\mathrm{nom}}=1963$~pb; shifts for each mask definition.}
\label{tab:phisym}
\begin{tabular}{lr}
\toprule
Mask & Shift [\%] \\
\midrule
Preselection & $+2.01$ \\
Common       & $+2.26$ \\
Tight        & $+1.37$ \\
Logical OR   & $+2.28$ \\
\midrule
Per-bin envelope & $\pm1.57$ \\
\bottomrule
\end{tabular}
\end{table}

\section{\textsc{jetphox} PDF dependence}
\label{app:jetphox}
The NLO theory reference used for comparison with the measured cross section is
\textsc{jetphox} (v1.3.1.4) with the CT14nlo parton distribution functions and
the BFG~II photon fragmentation functions. \textsc{jetphox} computes the
isolated-photon cross section at next-to-leading order in $\alpha_s$, including
both the direct contribution (the photon produced in the hard matrix element) and
the fragmentation contribution (the photon produced in the collinear
fragmentation of a final-state parton), with the fragmentation piece evolved
through the BFG fragmentation functions. To quantify the PDF-driven spread of the
theory prediction --- and in particular its sensitivity to the proton gluon
density, which the Compton subprocess $qg\to q\gamma$ makes the dominant initial
state at RHIC --- the production is rerun with two alternative PDF sets, the LO
CT14lo set and the NLO NNPDF3.1 set ($\alpha_s(M_Z)=0.118$), using identical
kinematic and fragmentation settings.

\paragraph{Method.} All three productions share the same configuration:
$\sqrts=200$~GeV $p+p$, photon rapidity $|y|<0.7$, direct$+$fragmentation
components with BFG~II, a fixed-cone truth isolation with $R=0.3$ and
$\Eiso^{\mathrm{truth}}<4$~GeV matching the analysis definition, and
renormalization/factorization/fragmentation scales
$\mu_R=\mu_F=\mu_f=c_m\cdot p_{\mathrm{T}}$ with $c_m\in\{0.5,1.0,2.0\}$ to map the
scale uncertainty. Per scale point, $100$ segments of $10^6$ events each are
generated, totaling $10^8$ events. The full PDF error set is computed for every
point ($56$ CT14 Hessian eigenvector pairs for CT14lo/CT14nlo, $100$ Monte Carlo
replicas for NNPDF3.1). The three productions are statistically independent: only
the LHAPDF identifier differs and the random-number seeds are offset so that no
event streams are shared. Pairing the LO CT14lo set with the NLO matrix element is
a deliberate LO-in-NLO inconsistency --- the missing NLO gluon evolution is
absorbed into an effective LO fit that is not strictly valid at all $x$ when
convolved with the NLO hard cross section --- and is shown here only as a
methodological probe of the LO-versus-NLO PDF effect, not as a physical
prediction.

\paragraph{Results.} Table~\ref{tab:jetphox} lists the central-scale
($\mu=1.0\cdot p_{\mathrm{T}}$) per-bin cross section for the three PDF sets,
with the ratios to CT14lo. Moving from the LO CT14lo set to an NLO set shifts the
prediction monotonically in $\ETg$: at low $\ETg$ ($8$--$14$~GeV) the NLO PDFs lie
$3$--$11\%$ \emph{above} CT14lo; at mid $\ETg$ ($14$--$26$~GeV) they lie
$4$--$25\%$ below; and at the highest bins ($26$--$36$~GeV) they lie $27$--$50\%$
below. The crossover near $\ETg\approx14$~GeV reflects the well-known
LO-versus-NLO gluon difference: at the moderate $x\sim2\ETg/\sqrts\sim0.1$
accessed at low $\ETg$ the LO gluon undershoots because NLO $K$-factors are
absorbed into the LO fit, whereas at higher $x\gtrsim0.2$ the LO fit
over-compensates to match jet data with the LO matrix element, an excess that the
NLO matrix element times NLO PDF correctly suppresses. The two NLO PDFs
(CT14nlo and NNPDF3.1) agree with each other to $\pm3\%$ at low $\ETg$ and grow
apart to $\pm20\%$ in the highest bin, where the PDF-family spread becomes
comparable to the renormalization-scale envelope ($\pm40\%$ from
$\mu=0.5/1.0/2.0$). This confirms that the LO-versus-NLO PDF effect, rather than
the specific NLO fit family, is the dominant PDF-driven theory uncertainty.

\begin{longtable}{lrrrrr}
\caption{\textsc{jetphox} cross section [pb] per $\ETg$ bin for three PDF sets
and ratios to CT14lo.}\label{tab:jetphox}\\
\toprule
$\ETg$ [GeV] & CT14lo & CT14nlo & NNPDF3.1 & nlo/lo & nnpdf/lo \\
\midrule
\endfirsthead
\toprule
$\ETg$ [GeV] & CT14lo & CT14nlo & NNPDF3.1 & nlo/lo & nnpdf/lo \\
\midrule
\endhead
\bottomrule
\endfoot
8--10   & 1519  & 1673  & 1690  & 1.101 & 1.113 \\
10--12  & 539.9 & 571.2 & 575.4 & 1.058 & 1.066 \\
12--14  & 218.5 & 222.2 & 225.1 & 1.017 & 1.030 \\
14--16  & 98.4  & 96.2  & 94.4  & 0.978 & 0.959 \\
16--18  & 47.0  & 44.0  & 41.9  & 0.937 & 0.891 \\
18--20  & 23.5  & 20.8  & 19.1  & 0.887 & 0.815 \\
20--22  & 9.82  & 8.60  & 8.01  & 0.876 & 0.816 \\
22--24  & 6.54  & 5.53  & 4.95  & 0.846 & 0.756 \\
24--26  & 3.66  & 2.93  & 2.73  & 0.799 & 0.745 \\
26--28  & 2.10  & 1.41  & 1.37  & 0.669 & 0.651 \\
28--32  & 0.940 & 0.608 & 0.579 & 0.647 & 0.616 \\
32--36  & 0.323 & 0.236 & 0.165 & 0.731 & 0.510 \\
\midrule
Integrated 8--36 & 4942 & 5294 & 5330 & $+7.1\%$ & $+7.9\%$ \\
\end{longtable}

\section{Back-to-back jet NPB cross-check}
\label{app:b2b}
This appendix provides an independent, physics-based validation of the
non-physics-background BDT that does not rely on the timing label used to train
it. The NPB BDT is trained against beam-induced streak clusters that are not
associated with a real $p+p$ collision and therefore carry no recoil jet. By
contrast, a genuine direct or fragmentation photon (and a high-$\ET$ $\pi^0$
candidate) is produced in a $2\to2$ hard scattering and is accompanied by a
balancing parton in the opposite hemisphere. Requiring a back-to-back recoil jet
(``b2b jet'') therefore selects a sample that is physically depleted of NPB
clusters, independently of any shower-shape-based veto. Comparing the unfolded
cross section in this b2b-enriched sample with and without the NPB BDT cut is thus
a clean closure test: if the BDT is unbiased on real photons, the two should agree
to within statistics.

\paragraph{Method.} Four variants are produced on top of the analysis pipeline.
\texttt{b2bjet\_pt5} adds a recoil-jet requirement with $p_{\mathrm{T}}^{\mathrm{jet}}\ge5$~GeV
back-to-back in $\phi$ ($|\Delta\phi|>3\pi/4$), keeping the nominal NPB cut
$\mathrm{NPB}>0.5$. \texttt{b2bjet\_pt5\_npb0} is the same but with the NPB cut
switched off. \texttt{b2bjet\_pt7} and \texttt{b2bjet\_pt7\_npb0} repeat these with
the higher threshold $p_{\mathrm{T}}^{\mathrm{jet}}\ge7$~GeV. Each variant runs the
full data$+$MC pipeline (purity, unfolding, efficiency, luminosity normalization)
at the same $\mathcal{L}=64.37~\mathrm{pb}^{-1}$ as the nominal, and the resulting
differential cross section is compared bin-by-bin to the nominal selection.

\paragraph{Results.} Two conclusions follow (original Figs.~69 and~70). First, the
NPB-on and NPB-off curves of the same recoil-jet threshold overlay tightly across
the analysis range: bin-by-bin,
$|\sigma_{\mathrm{pt5}}/\sigma_{\mathrm{nom}}-\sigma_{\mathrm{pt5,npb0}}/\sigma_{\mathrm{nom}}|\le1.0\%$
and the corresponding $p_{\mathrm{T}}^{\mathrm{jet}}\ge7$~GeV difference is
$\le0.5\%$, with the largest values in the lowest bin. The NPB cut therefore does
not change the answer once a genuine recoil jet is required, exactly the behavior
expected if the BDT is unbiased on real photons. Second, the b2b-enriched cross
section sits $6$--$10\%$ below the nominal in the lowest bin ($12$--$14$~GeV) and
recovers to within $\pm10\%$ of nominal for $\ETg\ge14$~GeV. The deficit is larger
for the $p_{\mathrm{T}}^{\mathrm{jet}}\ge7$~GeV variant ($10.0\%$) than for the
$\ge5$~GeV variant ($6.9\%$), which is the ordering expected from recoil-jet
\emph{acceptance}: at low $\ETg$ the recoil parton is soft and hadronizes into a
low-$\pt$ jet that the b2b cut sometimes rejects, more often at the higher
threshold. This interpretation is confirmed directly by the truth-level recoil-jet
probability (original Fig.~70), which rises from about $36\%$ ($24\%$) at
$\ETg=10$--$12$~GeV for $p_{\mathrm{T}}^{\mathrm{jet}}\ge5$ ($7$)~GeV to about
$75\%$ ($71\%$) at $\ETg=32$--$36$~GeV. The low-$\ETg$ deficit is therefore a
recoil-jet-acceptance effect inherited from the \textsc{pythia-8} fragmentation
modeling, not an NPB-cut bias, and is not propagated as a systematic on the
nominal cross section.

\section{PHENIX comparison corrections}
\label{app:phenix}
The PHENIX measurement of inclusive direct photons (Ref.~[1]) is reported in the
rapidity acceptance $|\eta|<0.25$ as a representative cross section per $\ETg$ bin
derived from the invariant cross section $E\,d^3\sigma/dp^3$. To overlay it on the
present isolated-photon measurement at $|\eta^{\gamma}|<0.7$ on an equal footing
(original Fig.~55), two corrections are applied to the published points: a
bin-width recovery from a smooth fit, and a pseudorapidity-density rescaling
derived from prompt-photon truth simulation. Both corrections preserve the
published central values and their fractional statistical uncertainties; they only
remove geometric and binning differences between the two experiments.

\paragraph{Bin-width recovery.} The PHENIX invariant cross section is first
converted to the differential form $d^2\sigma/(d\eta\,d\ETg)$ through the Jacobian
$2\pi\,\ETg$, and the result is fit over the published $5$--$26$~GeV range with a
smooth modified power law
\begin{equation}
  f(\ETg) = A\left(1 + \frac{(\ETg)^{2}}{b}\right)^{c},
  \qquad A=4.53\times10^{5},\; b=6.34,\; c=-3.602,
\end{equation}
with $\chi^2/\mathrm{ndf}=28.07/15$ and per-bin pulls bounded by $|2.5|$. Because
PHENIX quotes each point at the bin center after its own bin-shift correction,
converting to a strict bin average requires only a smooth multiplicative shape
factor read off the fit,
\begin{equation}
  y_i^{\mathrm{bin\text{-}avg}} = y_i^{\mathrm{pub}}\,
  \frac{\langle f\rangle_i}{f((\ETg)^c_i)},
  \qquad
  \langle f\rangle_{\mathrm{bin}}
  = \frac{1}{\Delta\ETg}\int_{\mathrm{bin}} f(\ETg)\,d\ETg,
\end{equation}
where $(\ETg)^c_i$ is the bin center. Only the few-percent shape distortion of a
steeply falling spectrum within finite-width bins is removed; the published value
and its statistical scatter are preserved.

\paragraph{Pseudorapidity-density rescaling.} The two experiments cover different
rapidity intervals, so the spectra are placed on a common acceptance via the
density ratio
\begin{equation}
  R(\ETg) = \frac{(dN/d\eta)|_{|\eta|<0.25}}{(dN/d\eta)|_{|\eta|<0.7}},
\end{equation}
evaluated from prompt-photon truth simulation (\textsc{pythia-8},
direct$+$fragmentation, no isolation requirement) on the analysis truth-$\pt$
grid (original Fig.~72). The ratio is consistent with unity to within $0.5\%$
below $\ETg\approx14$~GeV and rises monotonically to $R\approx1.22$ at
$\ETg=36$--$45$~GeV, reflecting the narrowing of the rapidity distribution of
$2\to2$ prompt-photon production at the higher Bjorken-$x$ accessed at high
$\ETg$. Each bin-averaged PHENIX value is divided by $R((\ETg)^c_i)$ to produce
the $|\eta|<0.7$-equivalent spectrum. Imposing the truth isolation
$\Eiso^{\mathrm{truth}}<4$~GeV on the prompt-photon sample shifts $R(\ETg)$ by
less than $0.1\%$ in every bin, so the same factor applies to PHENIX's inclusive
fiducial within the precision of either measurement.

\section{MBD photon/jet cross-check}
\label{app:mbd}
The MBD-vertex efficiency --- the probability that the minimum-bias detector
reconstructs a valid event vertex, $z_{\mathrm{reco}}\ne-9999$ --- enters the
cross section as one of the four efficiency factors of Sec.~\ref{sec:analysis}.
Because this efficiency depends on the forward particle activity that fires the
MBD, it could in principle differ between the photon-triggered topology and a
generic jet topology. This appendix verifies that it does not, which justifies
transferring the MBD systematic between channels (Sec.~\ref{sec:syst}, item~5.5).

The MBD-vertex-found fraction is computed in \textsc{pythia-8} as a function of
the leading truth $\pt$, separately for the photon channel (the
\texttt{photon\,5/10/20} samples, using the leading truth-photon $\pt$) and for
the jet channel (the \texttt{jet\,5/8/12/20/30/40} samples, using the leading
$R=0.4$ truth-jet $\pt$). Per-sample truth-$\pt$ windows
(Table~\ref{tab:ptwin}) prevent double-counting across sample boundaries, and each
event carries the weight $\sigma_i/N_i^{\mathrm{gen}}$. The two channels overlay
to within $1\%$ across the full reported range (original Fig.~73) and fall
monotonically from about $0.80$ at $p_{\mathrm{T}}^{\mathrm{lead}}=5$~GeV to about
$0.66$ at $35$~GeV. The decline reflects the reduced forward activity at higher
Bjorken-$x$: as the hard scale rises, more of the collision energy is carried into
the central region and less into the forward MBD acceptance, lowering the
coincidence probability. The channel-level agreement demonstrates that the
MBD-vertex efficiency is governed by the underlying-event and beam-remnant
activity rather than by the hard-process flavor, which is what permits the
data-driven MBD efficiency (derived in the PPG-10 luminosity note) to be applied
to the photon measurement.

\section{Double-interaction (pileup) modeling}
\label{app:di}
At the luminosities reached in Run~24 a non-negligible fraction of bunch crossings
contain more than one $p+p$ collision. A second, overlapping collision affects the
photon measurement in two distinct ways: it adds soft energy that can enter the
isolation cone, and --- more importantly --- it biases the reconstructed event
vertex, which in turn shifts the inferred cluster pseudorapidity and distorts the
shower-shape inputs to the photon-ID BDT. This appendix develops the full
double-interaction (DI) model used throughout the analysis. Rather than applying a
factorized after-the-fact correction, the analysis builds dedicated full-\textsc{geant}
DI samples and \emph{blends} them with the single-interaction (SI) samples at the
physically determined cluster fraction, so that the efficiency, the response
matrix, and the background template all inherit the pileup effects consistently.

\paragraph{DI samples.} Dedicated \texttt{photon\{5,10,20\}\_double} and
\texttt{jet\{8,12,20,30,40\}\_double} samples are produced by full simulation of
two overlapping $p+p$ collisions: one hard-scatter event drawn from the same
\textsc{pythia-8} generator as the corresponding SI sample, plus one minimum-bias
\textsc{pythia-8} (Detroit-tune) event overlaid as pileup. The two collisions are
placed at independently drawn primary vertices from the generator beam profile.
These DI samples are mixed with the nominal SI samples at the cluster-weighted
fractions of Table~\ref{tab:di}, providing both the DI efficiency study and the
shower-shape data--MC validation below.

\paragraph{Pileup fractions.} The mean number of interactions per crossing is
obtained from the Poisson-corrected MBD coincidence rate,
\begin{equation}
  \mu = \frac{\sum_{\mathrm{runs}} N_{\mathrm{MBD,corr}}}{r_B\,\sum_{\mathrm{runs}} \Delta t},
  \qquad r_B = 111\times78\,200 = 8.68\times10^{6}~\mathrm{Hz},
\end{equation}
where $r_B$ is the RHIC bunch-crossing rate. The event-level interaction fractions
follow a zero-truncated Poisson (one collision is guaranteed by the trigger),
\begin{equation}
  f_1 = \frac{\mu e^{-\mu}}{1-e^{-\mu}},\quad
  f_2 = \frac{(\mu^2/2)e^{-\mu}}{1-e^{-\mu}},\quad
  f_{\ge3} = 1-f_1-f_2,
\end{equation}
computed per run and luminosity-weighted. Because multi-interaction events carry a
higher cluster multiplicity, the relevant quantities for the cluster-level blend
are the cluster-weighted fractions
\begin{equation}
  w_{\mathrm{single}} = \frac{f_1}{\sum_k k\,f_k},\qquad
  w_{\mathrm{double}} = \frac{2 f_2 + 3 f_{\ge3}}{\sum_k k\,f_k},
\end{equation}
with triple-and-higher interactions folded into the double bin (they are rare
enough that the two-collision sample models them adequately). The resulting
fractions and weights are given in Table~\ref{tab:di}; the cluster-weighted
$w_{\mathrm{double}}$ values ($0.224$ at $0$~mrad, $0.079$ at $1.5$~mrad) are the
numbers used throughout the analysis.

\begin{table}[h]
\centering
\caption{Double-interaction parameters by crossing-angle configuration.}
\label{tab:di}
\begin{tabular}{lcc}
\toprule
Quantity & 0~mrad & 1.5~mrad \\
\midrule
$\mu_{\mathrm{corr}}$ & 0.24 & 0.08 \\
$f_1$       & 0.879 & 0.960 \\
$f_2$       & 0.111 & 0.039 \\
$w_{\mathrm{single}}$ & 0.776 & 0.921 \\
$w_{\mathrm{double}}$ & 0.224 & 0.079 \\
\bottomrule
\end{tabular}
\end{table}

\paragraph{Vertex-shift mechanism.} In a DI event the MBD reconstructs a single
vertex whose best estimator is the average of the two collision points,
\begin{equation}
  z_{\mathrm{reco}} \approx \frac{z_1 + z_2}{2}.
\end{equation}
The two true vertices are drawn from the luminous region (data) or the generator
profile (MC), so the displacement between the primary truth vertex and the
reconstructed vertex, $\Delta z\equiv z_{\mathrm{reco}}-z_{\mathrm{truth}}$, can
reach several tens of cm. A cluster reconstructed at the EMCal barrel radius
$R_{\mathrm{EMCal}}=93.5$~cm then acquires a pseudorapidity shift that follows
directly from the geometry,
\begin{equation}
  \Delta\eta \approx -\frac{\Delta z}{R_{\mathrm{EMCal}}\cosh\eta},
\end{equation}
reducing at $\eta=0$ to $|\Delta\eta|\approx|\Delta z|/93.5~\mathrm{cm}$. This
$\eta$ shift is the root cause of the shower-shape distortion that drives the
photon-ID degradation below. The RMS of $\Delta z$ grows from $7.6$~cm in SI MC to
$34.5$~cm in DI MC (Table~\ref{tab:dz}).

A per-event check confirms that the second collision does not add genuine
prompt-photon content: the mean number of truth photons per event is $0.493$
(single) versus $0.495$ (double), and the mean truth isolation energy is
$0.17$ versus $0.26$~GeV, both far below the $4$~GeV fiducial threshold. The DI
effect is therefore a reconstruction effect, not a change in the underlying
signal.

\paragraph{Per-stage efficiency.} Truth-to-reco association uses the
\textsc{geant} track-ID from the \texttt{CaloRawClusterEval} module as the
unambiguous ground truth, with no additional geometric requirement (so that a
fully reconstructed but $\eta$-shifted cluster is still counted as found). The
factorized chain
$\varepsilon_{\mathrm{reco}}\to\varepsilon_{\mathrm{iso}}\to\varepsilon_{\mathrm{id}}\to\varepsilon_{\mathrm{all}}$
then isolates the stage at which each pileup effect acts; the integrated values
for SI, DI, and the $0$~mrad mixed sample are given in Table~\ref{tab:dieff}.
The reconstruction stage loses $20$ percentage points (pp), of which about
$5.5$~pp is genuine cluster loss (no track-ID match: $7.9\%$ in DI versus $2.4\%$
in SI) and about $15$~pp is common-cut failure from shifted kinematics. The
isolation stage is barely affected ($-3$~pp), since the extra cone energy is small
compared with the $4$~GeV threshold. The photon-ID stage loses $13$~pp, entirely
driven by the vertex-displaced shower-shape inputs feeding the $\ET$-dependent BDT
threshold: the vertex-shifted population has a BDT pass rate of only $\sim60\%$
versus $82.7\%$ for the unshifted population. The integrated DI-induced efficiency
reduction is about $9\%$ at $0$~mrad and about $3\%$ at $1.5$~mrad, scaling with
the cluster-weighted pileup fraction.

These effects are absorbed directly into the nominal pipeline rather than applied
as a separate correction: the efficiency, the response matrix, and the
inclusive-jet background template are all built from the SI$+$DI mixed ensemble,
with each event carrying the weight
$w = \sigma_{\mathrm{eff}}\cdot f_{\mathrm{mix}}\cdot\mathrm{TVR}(z_h,z_{mb})$,
where $f_{\mathrm{mix}}=1-f_{\mathrm{DI}}$ for SI samples and $f_{\mathrm{DI}}$ for
DI samples, and $\mathrm{TVR}$ is the truth-vertex reweight of
Appendix~\ref{app:tvr}. The single free parameter $f_{\mathrm{DI}}$ is varied to
obtain the systematic of item~5.8.

\begin{table}[h]
\centering
\caption{RMS of $\Delta z$ for single and double interactions [cm].}
\label{tab:dz}
\begin{tabular}{lc}
\toprule
Class & RMS $\Delta z$ [cm] \\
\midrule
Single & 7.6 \\
Double & 34.5 \\
\bottomrule
\end{tabular}
\end{table}

\begin{table}[h]
\centering
\caption{Per-stage efficiency [\%] for single, double and mixed events, and the
net effect (percentage points).}
\label{tab:dieff}
\begin{tabular}{lcccc}
\toprule
Stage & Single & Double & Mixed & Effect \\
\midrule
$\varepsilon_{\mathrm{reco}}$ & 93.5 & 73.1 & 88.4 & $-20$~pp \\
$\varepsilon_{\mathrm{iso}}$  & 74.2 & 71.0 & 73.5 & $-3$~pp  \\
$\varepsilon_{\mathrm{id}}$   & 82.7 & 69.4 & 80.0 & $-13$~pp \\
$\varepsilon_{\mathrm{all}}$  & 57.4 & 36.0 & 52.0 & $-21$~pp \\
\bottomrule
\end{tabular}
\end{table}

\paragraph{Shower-shape data--MC validation.} Before fitting the blend fraction,
the model is validated at the sample level. The most pileup-sensitive observable
is the center-of-gravity width $w_{\mathrm{CoG}\eta}$, which is broadened directly
by the vertex shift via the $\Delta\eta$ relation above; the photon-ID BDT score
amplifies the residual through its nonlinear discriminant, while the
energy-spread-driven observables ($f_1$, $E_{1\times1}/E_{3\times3}$) are far less
DI-sensitive. The original Figs.~75--76 overlay pure SI MC, pure DI MC, and data
from both crossing-angle periods, all area-normalized at the preselection cut. The
$0$~mrad data systematically track the DI tail, while the $1.5$~mrad data track
closer to the SI MC --- exactly as expected from their $22.4\%$ versus $7.9\%$
pileup fractions --- providing a direct sample-level confirmation of the blending
picture. Representative data$-$inclusive-MC $\chi^2/\mathrm{ndf}$ values for the
two leading discriminants at the lowest and highest $\ETg$ bins are given in
Table~\ref{tab:dichi2}; the mid- and high-$\ETg$ bins improve dramatically once
the DI component is blended in, whereas the lowest bin is dominated by a residual
DI--MC mismodeling discussed below.

\paragraph{Data-driven blend-fraction scan.} The nominal blending uses the
analytic cluster-weighted fractions $f_{\mathrm{DI}}=0.224$ ($0$~mrad) and
$0.079$ ($1.5$~mrad). The systematic uncertainty on $f_{\mathrm{DI}}$ is sized by
treating the fraction as a free parameter and minimizing the $\chi^2$ between the
data and an inclusive MC rebuilt by shape interpolation,
\begin{equation}
  \mathrm{MC}_{\mathrm{incl}}(f) = (1-f)\,T_{\mathrm{SI}} + f\,T_{\mathrm{DI}}.
\end{equation}
The pure SI and DI templates $T_{\mathrm{SI}},T_{\mathrm{DI}}$ are recovered by
``de-baking'' the analytic weights out of the combined hadds,
$T_{\mathrm{SI}}=\mathrm{SI}_{\mathrm{total}}/(1-f_{\mathrm{DI}}^{\mathrm{anal}})$
and
$T_{\mathrm{DI}}=\mathrm{DI}_{\mathrm{total}}/f_{\mathrm{DI}}^{\mathrm{anal}}$,
and the round-trip identity
$(1-f)T_{\mathrm{SI}}+fT_{\mathrm{DI}}=\mathrm{SI}_{\mathrm{total}}+\mathrm{DI}_{\mathrm{total}}$
at $f=f_{\mathrm{DI}}^{\mathrm{anal}}$ is verified numerically. The scan is a
one-parameter shape-only $\chi^2$ against the data BDT-score distribution at the
preselection cut (NPB-clean, jet-dominated), unit-area-normalized in each $\ETg$
bin so that the absolute jet cross-section uncertainty does not leak into
$f_{\mathrm{DI}}$. The BDT score is chosen because it aggregates the eleven
shower-shape inputs through a discriminant already trained on the SI/DI
separation. Each crossing angle is scanned over $f\in[0.00,0.50]$ in steps of
$0.005$, summing residuals over five $\ETg$ cells ($5\times50=250$ bins). The
before/after $\chi^2/\mathrm{ndf}$ for the two variables and two $\ETg$ ranges is
given in Table~\ref{tab:dichi2}, and the scan minimum in Table~\ref{tab:discan}.

The $0$~mrad scan has a clean parabolic minimum at $f_{\mathrm{best}}=0.290$, above
the analytic anchor $0.224$; the $1.5$~mrad minimum sits at the lower scan edge
$f_{\mathrm{best}}=0$, below the analytic anchor $0.079$. The residual
$\chi^2/\mathrm{ndf}$ at the minimum is far from unity ($2.57$ and $3.20$), so the
$\Delta\chi^2=1$ prescription does not give a statistically meaningful interval;
moreover the two crossings pull $f_{\mathrm{best}}$ in \emph{opposite} directions
relative to their anchors, whereas a coherent DI-shape mismodeling would shift them
in the same fractional direction. This indicates that the $\chi^2$ surface is
dominated by data--MC tensions unrelated to the DI fraction (truth-vertex-reweight
closure, jet-fragmentation modeling, BDT calibration). Accordingly the central
prediction is kept at the analytic anchors, and the fit is used only to size the
envelope: a single variation per crossing angle is run at $f_{\mathrm{best}}$
($0.290$ at $0$~mrad, $0.000$ at $1.5$~mrad), and the absolute per-bin cross-section
deviation is taken as the symmetric two-sided systematic of item~5.8.

\begin{table}[h]
\centering
\caption{$\chi^2/\mathrm{ndf}$ before $\to$ after DI blending, for two variables
and two $\ETg$ ranges, per crossing-angle configuration.}
\label{tab:dichi2}
\begin{tabular}{llcc}
\toprule
Variable & $\ETg$ range & 0~mrad & 1.5~mrad \\
\midrule
$w_{\mathrm{CoG}\eta}$ & 10--14 & $47.40\to28.96$ & $9.72\to42.31$ \\
$w_{\mathrm{CoG}\eta}$ & 22--28 & $15.58\to0.75$  & $21.10\to1.31$ \\
BDT                    & 10--14 & $6.85\to20.30$  & $5.98\to33.08$ \\
BDT                    & 22--28 & $9.62\to1.75$   & $1.89\to1.13$  \\
\bottomrule
\end{tabular}
\end{table}

\begin{table}[h]
\centering
\caption{DI blend-fraction scan: analysis value, best-fit value, before/after
$\chi^2/\mathrm{ndf}$ and $\Delta\chi^2$.}
\label{tab:discan}
\begin{tabular}{lcc}
\toprule
Quantity & 0~mrad & 1.5~mrad \\
\midrule
$f_{\mathrm{anal}}$ & 0.224 & 0.079 \\
$f_{\mathrm{best}}$ & 0.290 & 0.000 \\
$\chi^2/\mathrm{ndf}$ (before) & 2.57 & 3.20 \\
$\chi^2/\mathrm{ndf}$ (after)  & 2.84 & 3.67 \\
$\Delta\chi^2$ & 68.4 & 116.6 \\
\bottomrule
\end{tabular}
\end{table}

\paragraph{Cluster size under DI.} The EMCal cluster size (the owned-tower count
$n_{\mathrm{owned}}$ and the bounding-box widths $\Delta i_\eta,\Delta i_\phi$) is
the basic shape property underlying every shower-shape discriminant. The
DI-induced shifts are mild for photons ($\Delta\langle n_{\mathrm{owned}}\rangle
\sim0.2$ towers across the analysis range) but several times larger for jets,
because soft neighboring towers from the second collision can be pushed above the
$70$~MeV clustering threshold in a busy jet environment. After the tight photon-ID
selection, data, jet MC, and photon MC all converge to a signal-like cluster size
within $\pm5\%$, so the residual DI effect on the ABCD background template is most
visible at the preselection stage and is reduced to below about $1\%$ after the
tight cut.

\paragraph{Energy response under DI.} The photon energy response
$R=E_{\mathrm{T}}^{\mathrm{reco}}/p_{\mathrm{T}}^{\mathrm{truth}}$ is measured in
bins of truth $\pt$ across the SI, DI, and mixed samples by fitting a
double-sided Crystal-Ball function to extract the peak $\mu$ and resolution
$\sigma$. At the common-cut level the DI events shift the peak by $\lesssim0.5\%$
and broaden the resolution by $\lesssim10\%$ at $\pt\sim20$~GeV, with the mixed
curves sitting between SI and DI at the appropriate fractions. This shift is well
inside the energy-scale systematic budget and is further reduced by the
truth-vertex reweight, so the leading DI sensitivity remains on the efficiency and
the background template rather than on the energy scale.

\paragraph{DI reco--truth kernel.} A structural feature limits any truth-only
vertex correction for DI events. In \texttt{jet12\_double} MC at $1.5$~mrad the
conditional reco kernel $P(z_r\,|\,z_h=0)$, with the minimum-bias vertex $z_{mb}$
marginalized, has an RMS of $\sim34$~cm --- wider than the $\sim22$~cm data target.
Because the reconstructed vertex of a DI event is the average of two
independently drawn vertices, even forcing the primary truth vertex to $z_h=0$
cannot narrow the reco distribution below this kernel width. This is the
geometric reason a reweight on the primary truth vertex $z_h$ alone is
insufficient for DI events, and it motivates the two-vertex
$w(z_h)\,w(z_{mb})$ factorization adopted in the next appendix.

\section{Truth-vertex reweighting}
\label{app:tvr}
The longitudinal extent of the luminous region differs between the two
crossing-angle periods ($\sigma_{\mathrm{lum}}\approx54$~cm at $0$~mrad,
$\approx22$~cm at $1.5$~mrad) and from the common generator beam profile
($\sigma_{\mathrm{gen}}\approx65$~cm). Since the vertex-cut acceptance and the
vertex-driven shower shapes both depend on the $z$ distribution, the MC
truth-vertex distribution must be reweighted to match data. A naive reco-level
weight $f(z_r)$ is insufficient for DI events: $f(z_r)$ depends only on the
symmetric coordinate $\bar z=(z_h+z_{mb})/2$ and is flat in the antisymmetric
coordinate $\varepsilon=(z_h-z_{mb})/2$, so the antisymmetric component is
unconstrained by reco closure alone and the primary truth vertex retains a width
$\sigma(z_h)=\sqrt{\tfrac12\sigma_{\mathrm{lum}}^2+\tfrac12\sigma_{\mathrm{MC}}^2}\simeq47$~cm
at $1.5$~mrad. Left uncorrected, this biases the DI truth fiducial acceptance low
by about $19$~pp, propagating to a $+1.4\%$ bias on the cross section at $1.5$~mrad
(and $\sim0.3\%$ at $0$~mrad). The reweight is therefore factorized as $w(z_h)$ for
SI events and $w(z_h)\,w(z_{mb})$ for DI events, giving the iteration access to the
antisymmetric coordinate.

\paragraph{Method.} The weight function $w(z)$ is derived directly from the data
reco distribution with no closed-form Gaussian assumption. Under a candidate
weight, the mixed-MC reco distribution is built event-by-event,
\begin{equation}
  M(z_r) = f_s\!\!\sum_{e\in T_s}\! w(z^e_h)\,\delta(z_r-z^e_r)
         + f_d\!\!\sum_{e\in T_d}\! w(z^e_h)\,w(z^e_{mb})\,\delta(z_r-z^e_r),
\end{equation}
where $f_d$ is the cluster-weighted DI fraction and $f_s=1-f_d$; for DI events the
weight enters once per truth vertex. Both periods use $80$ bins over
$[-200,200]$~cm, with the closure window auto-derived as the $99$th percentile of
$|z_r^{\mathrm{data}}|$ ($z_{\mathrm{cut}}=83$~cm at $1.5$~mrad, $144$~cm at
$0$~mrad). The seed is
$w_0(z)=\mathcal{N}(0,\sigma_{\mathrm{seed}}^2)/\mathcal{N}(0,\sigma_{\mathrm{gen}}^2)$
with $\sigma_{\mathrm{gen}}=64.997\pm0.015$~cm from the MC truth vertex and
$\sigma_{\mathrm{seed}}=\mathrm{std}(z_r^{\mathrm{data}})$ ($21.48$ / $54.32$~cm);
$\sigma_{\mathrm{seed}}$ needs no physics input --- it is simply the empirical
width of the distribution being matched. The update is Lucy--Richardson-like,
multiplicative, and damped,
\begin{align}
  z_{\mathrm{reco}}^{\mathrm{model}} &= \mathcal{F}(z_{\mathrm{truth}}), \\
  w_{i+1}(z) &= w_i(z) \cdot \big[\mathrm{smooth}(r_i(z))\big]^{\alpha},
  \qquad \alpha=0.5,
\end{align}
with $r_i(z_r)=D(z_r)/M_i(z_r)$ the data/MC residual applied at the matching truth
bin. After each step $w$ is renormalized so $\langle w(z_h)\rangle=1$ over the SI
truth distribution (preserving the physics DI fraction); one-bin Gaussian
smoothing suppresses oscillation, the update factor is clipped to $[0.5,2]$ per
iteration, and boundary bins are edge-clamped. Convergence is declared at
$\max|r_i(z)-1|<0.01$ inside the window, with a $30$-iteration cap. The converged
parameters are listed in Table~\ref{tab:tvr}.

\begin{table}[h]
\centering
\caption{Truth-vertex-reweighting parameters by crossing-angle configuration.}
\label{tab:tvr}
\begin{tabular}{lcc}
\toprule
Quantity & 0~mrad & 1.5~mrad \\
\midrule
Luminosity [pb$^{-1}$] & 47.2076 & 17.1642 \\
$f_{\mathrm{double}}$ & 0.224 & 0.079 \\
$\sigma_{\mathrm{seed}}$ [cm] & 54.32 & 21.48 \\
$\sigma_{\mathrm{gen}}$ [cm] & \multicolumn{2}{c}{$64.997 \pm 0.015$} \\
Final $\chi^2/\mathrm{ndf}$ & 2.78 (ndf 57) & 13.7 (ndf 33) \\
$\max|r-1|$ & 0.008 & 0.111 \\
\bottomrule
\end{tabular}
\end{table}

\paragraph{Per-period results.} Both periods are fit on the
\texttt{photon10}/\texttt{photon10\_double} mix; the headline closure numbers are
in Table~\ref{tab:tvr}. The two cases behave very differently. At $0$~mrad the
luminous region ($\sigma_{\mathrm{lum}}\approx54$~cm) is already close to the
generator profile and the data reco shape is nearly Gaussian, so the seed alone
gives $\max|r-1|\approx8\%$ and the iteration converges cleanly in $8$ steps to
$\max|r-1|\approx0.8\%$ and $\chi^2/\mathrm{ndf}\approx2.78$; the recovered $w(z)$
is smooth and broadly Gaussian with a small central enhancement. At $1.5$~mrad the
narrow luminous region must be extracted from a roughly three-times-wider source,
so $w(0)\approx8$, and the data reco distribution is non-Gaussian (a narrow core
plus a mid-tail shoulder near $|z_r|\approx45$--$60$~cm). The iteration makes
dramatic progress in the first $\sim5$ steps (the in-window $\chi^2$ falls by a
factor $\sim3100$) and then enters a slow plateau driven by bins at the
closure-window edge $|z_r|\approx83$~cm, where the wide DI reco--truth kernel
(previous appendix) makes the residual insensitive to $w(z)$. After $30$ iterations
the bulk $|z_r|<30$~cm is closed to $1$--$2\%$ per bin while the edge residual
stalls at $\max|r-1|\approx0.11$. Inside the analysis fiducial $|z|<60$~cm the
per-bin closure is at the $1$--$2\%$ level for both periods.

A single-Gaussian ansatz cannot capture the $1.5$~mrad shape: a parametric
$\sigma_{\mathrm{tgt}}$ scan bottoms out at $\sigma_{\mathrm{tgt}}\approx15.7$~cm
with a $\chi^2/\mathrm{ndf}$ two orders of magnitude worse than the iterative
result, whereas at $0$~mrad the parametric and iterative solutions nearly coincide
because the data is already Gaussian. The recovered $1.5$~mrad $w(z)$ has three
distinguishing features that no symmetric Gaussian can represent: a strong central
peak ($w(0)\approx8$) pulling near-IP doubles into the data-like core, a mid-$z$
dip suppressing the generator's natural $|z|\sim30$--$50$~cm population, and a
clear $\pm$ asymmetry in $|z|\approx40$--$60$~cm tracking the asymmetric data
mid-tail. The reweight is fit on the photon mix; extending it to the full sample
blend (including \texttt{jet12}) is expected to shift the closure metrics by about
$1\%$. The residual cross-section impact, estimated from the maximum-deviation
iteration, is $\mathcal{O}(1$--$2\%)$, concentrated at low $\ETg$ where the
vertex-cut acceptance is most sensitive; it is carried as a cross-check and
absorbed into the photon-ID response-driven systematic group rather than added in
quadrature.

\section{\textsc{herwig} cross-check}
\label{app:herwig}
The efficiency factors of Sec.~\ref{sec:analysis} are derived from
\textsc{pythia-8} MC and therefore inherit any mismodeling of the event topology
around the photon --- in particular the underlying-event and beam-remnant activity
that sets the isolation and MBD-vertex efficiencies. To bound this generator
dependence, the full $1.5$~mrad analysis pipeline is rerun with the nominal
\textsc{pythia-8} Detroit signal MC replaced by \textsc{herwig}~7 photon-jet
samples (\texttt{photon10\_herwig}, \texttt{photon20\_herwig}), keeping the data
input and all background MC unchanged. \textsc{herwig} uses an angular-ordered
parton shower and a cluster hadronization model, in contrast to \textsc{pythia}'s
dipole shower and Lund string fragmentation, so the comparison probes a genuinely
different modeling of the soft environment. The truth-$\pt$ prior is rederived in
the same way as the nominal (an unweighted single-iteration unfolding of the data
with the \textsc{herwig} response matrix, refit to a rational form), and the
unfolded comparison uses a common reweight target so that any residual difference
reflects the response matrix alone rather than the prior choice.

The differences are concentrated in the isolation and MBD-vertex efficiencies:
\textsc{herwig} photons sit in less-active mid-rapidity environments at fixed
truth $\ETg$, raising the isolation efficiency by $3$--$6\%$ and the MBD-vertex
efficiency by $4$--$7\%$. The photon-ID efficiency, which is governed by the
calorimeter shower physics rather than the event topology, agrees to within
$\pm1.5\%$ across the analysis range, and the unfolded yield (before efficiency
division) collapses to within $\pm4\%$ once the common reweight target is used. The
isolation-efficiency difference is propagated as the generator-driven systematic in
the efficiency group (Sec.~\ref{sec:syst}, item~5.4.2, $\pm5.3\%$). The
MBD-vertex difference is \emph{not} double-counted here, because the data-driven
MBD-coincidence efficiency from the PPG-10 luminosity note already covers the
generator-modeling uncertainty on the MBD-trigger response (item~5.5). The
unfolded-yield comparison is retained as a closure cross-check only.

\end{document}
